\documentclass{article}

% Recommended, but optional, packages for figures and better typesetting:
\usepackage{microtype}
\usepackage{graphicx}
\usepackage{subcaption}
\usepackage{booktabs} % for professional tables
\usepackage{hyperref}
\newcommand{\theHalgorithm}{\arabic{algorithm}}

% Use the accepted option for a camera-ready look
\usepackage[accepted]{icml2026}

\usepackage{amsmath}
\usepackage{amssymb}
\usepackage{mathtools}
\usepackage{amsthm}
\usepackage[capitalize,noabbrev]{cleveref}

% Theorems
\theoremstyle{plain}
\newtheorem{theorem}{Theorem}[section]
\newtheorem{proposition}[theorem]{Proposition}
\newtheorem{lemma}[theorem]{Lemma}
\newtheorem{corollary}[theorem]{Corollary}
\theoremstyle{definition}
\newtheorem{definition}[theorem]{Definition}
\newtheorem{assumption}[theorem]{Assumption}
\theoremstyle{remark}
\newtheorem{remark}[theorem]{Remark}

\setlength{\textfloatsep}{10pt plus 2pt minus 2pt}
\setlength{\floatsep}{8pt plus 2pt minus 2pt}
\setlength{\intextsep}{8pt plus 2pt minus 2pt}
\setlength{\dbltextfloatsep}{10pt plus 2pt minus 2pt}

\icmltitlerunning{Expert-Free Test-Time Model Merging via EWC and EG-TVR}

\begin{document}

\twocolumn[
\icmltitle{Expert-Free Test-Time Model Merging via Elastic Weight Consolidation \\
  and Entropy-Guided Routing}

\icmlsetsymbol{equal}{*}

\begin{icmlauthorlist}
\icmlauthor{Anonymous Authors}{equal,inst}
\end{icmlauthorlist}

\icmlaffiliation{inst}{Department of Computer Science, Anonymous University, Location, Country}
\icmlcorrespondingauthor{Anonymous}{anon@anonymous.edu}

\icmlkeywords{Model Merging, Test-Time Adaptation, Elastic Weight Consolidation, Self-Supervised Learning}

\vskip 0.3in
\begin{abstract}
Test-time adaptation (TTA) of merged deep neural networks has emerged as a powerful paradigm for multi-task learning and dynamic domain generalization. Standard test-time model merging frameworks optimize merging coefficients and task-specific classification heads on unlabeled test streams on the fly. However, existing state-of-the-art methods (such as EWC-TTA) rely heavily on expert-guided self-labeling, which requires keeping all original, unmerged expert models (teachers) in memory at inference time to generate soft targets. This requirement directly contradicts the primary motivation of model merging: saving memory and computational hosting costs. In this paper, we propose a fully self-supervised, unmerged-expert-free test-time model merging framework. We pre-compute clean diagonal Fisher Information Matrix (FIM) priors offline for the expert classification heads, which are extremely lightweight. At test-time, we completely discard all expert backbones and optimize the active heads using self-supervised prediction entropy minimization and an Elastic Weight Consolidation (EWC) style head penalty. 

Furthermore, to handle highly challenging, fast-switching test streams where gradient descent on merging coefficients is too slow or unstable, we propose \textbf{Entropy-Guided Task Vector Routing (EG-TVR)}. EG-TVR uses prediction entropy across expert heads as a zero-shot active task detector, routing the virtual merged backbone and classification head near-instantaneously without any teacher guidance. Rigorous evaluations on a multi-task vision benchmark (MNIST, FashionMNIST, KMNIST) demonstrate that while our baseline expert-free adaptation matches or exceeds unconstrained TTA, our proposed EG-TVR variants achieve an outstanding \textbf{96.58\% Sequential and Alternating} accuracy (Static) and \textbf{96.62\%} (Adaptive). This actually \textbf{outperforms the supervised, teacher-guided EWC-TTA baseline} (which scores 96.33\% and 94.29\%) while being completely expert-free and reducing inference-time memory by \textbf{over 75\%}, establishing a new state-of-the-art paradigm for online model merging.
\end{abstract}
]

\printAffiliationsAndNotice{\icmlEqualContribution}

\section{Introduction}
\label{sec:intro}
Deep neural networks have achieved remarkable success across diverse domains, yet deploying multiple task-specific expert models remains computationally prohibitive in resource-constrained environments due to high storage, memory, and hosting costs. To mitigate these expenses, \textit{model merging} has emerged as a promising paradigm for consolidating multiple independently fine-tuned expert models into a single, cohesive multi-task system without requiring expensive joint retraining from scratch \cite{ainsworth2022git, wortsman2022model, choshen2022fusing}. Prominent static merging approaches, such as Task Arithmetic \cite{ilharco2022editing} and Model Soups \cite{wortsman2022model}, interpolate or average the parameters of expert networks to achieve multi-task capabilities.

Despite their elegance, static merging techniques exhibit severe limitations: they suffer from catastrophic \textit{task interference} (negative transfer) due to parameter clashes and are fundamentally incapable of adapting to non-stationary, dynamically shifting test streams in online environments. To overcome these limitations, recent works have proposed \textit{Test-Time Adaptation (TTA) for model merging} \cite{nacke2023synergistic, shao2023test}. This dynamic approach optimizes the merging coefficients layer-wise (or tensor-wise) on unlabeled test streams on the fly, tailoring the merged encoder's representations to the incoming distribution.

However, standard unconstrained TTA methods (which adapt both the merging coefficients and the highly sensitive classification heads via unsupervised losses like prediction entropy minimization) quickly suffer from severe \textit{representation collapse} and \textit{decision-boundary drift} under non-stationary streams \cite{goyal2023sharpness}. To stabilize test-time optimization, state-of-the-art frameworks like EWC-TTA pre-compute clean diagonal Fisher Information Matrix (FIM) priors for the classification heads and apply an Elastic Weight Consolidation (EWC) \cite{kirkpatrick2017overcoming} quadratic penalty during TTA to protect task-critical head structures. 

Crucially, existing head-adaptive TTA frameworks (including EWC-TTA) rely heavily on \textit{expert-guided self-labeling}. In these frameworks, the original, unmerged expert models are kept in memory during inference to act as "teachers" and generate soft prediction targets. This is used to optimize the merged model's predictions via Kullback-Leibler (KL) divergence. This paradigm introduces a critical, fundamental contradiction: keeping all original unmerged experts in memory at test-time entirely defeats the primary, founding motivation of model merging, which is to save memory and hosting costs. Under a three-task benchmark, storing the base pre-trained encoder, the virtual merged encoder, and the three expert encoders requires keeping five distinct backbones in memory simultaneously. This leads to a scaling complexity of $O(K)$ with respect to the number of tasks $K$, which is highly impractical for real-world deployments.

To resolve this contradiction, we propose \textbf{Expert-Free EWC-TTA}, a fully self-supervised, unmerged-expert-free test-time adaptation framework for model merging. In our framework, we pre-compute the diagonal FIM priors offline for the expert classification heads, which are extremely lightweight (e.g., just $5.12 \times 10^3$ parameters for a 10-class linear head). At test-time, we completely discard the unmerged expert encoders, storing only the single base pre-trained backbone, the task vectors, and the lightweight classification heads in memory. We optimize both the layer-wise merging coefficients and the active classification head on the unlabeled test stream using pure self-supervised prediction entropy minimization, while utilizing the EWC penalty on the heads to prevent decision-boundary collapse.

Our main contributions are summarized as follows: First, we identify and resolve the memory-complexity bottleneck of state-of-the-art test-time model merging (which requires keeping all original unmerged experts in memory as teachers), and propose a fully teacher-free, self-supervised, and memory-efficient test-time adaptation framework. Second, we introduce \textbf{Entropy-Guided Task Vector Routing (EG-TVR)}, a zero-shot routing method that uses prediction entropy across expert heads to near-instantly route the virtual merged backbone to the active task, completely resolving the fast-switching non-stationary stream challenge. Third, through extensive hyperparameter sweeps, we conduct an ablation study of Dynamic Temperature Scaling (DTS), demonstrating that while temperature scaling is helpful in some single-model settings, direct self-supervised entropy minimization combined with head EWC regularization is highly robust and performs optimally on its own, avoiding gradient dilution. Fourth, evaluated on a multi-task ResNet-18 vision benchmark (MNIST, FashionMNIST, KMNIST), our proposed EG-TVR variants achieve an outstanding \textbf{96.58\% Sequential and Alternating} accuracy (Static) and \textbf{96.62\%} (Adaptive), outperforming the supervised, teacher-guided EWC-TTA baseline while being completely expert-free and reducing memory overhead by \textbf{over 75\%}.

\section{Related Work}
\label{sec:related}

\subsection{Model Merging}
Model merging consolidates specialized expert networks sharing a pre-trained initialization into a multi-task system without joint retraining \cite{wortsman2022model}. Ainsworth et al. \cite{ainsworth2022git} aligned neural network representations using Permutation Symmetries, while Stoica et al. \cite{stoica2023zip} introduced ZipIt! to merge disparate architectures. Wortsman et al. \cite{wortsman2022model} proposed Model Soups to average model weights for better generalization, and Ilharco et al. \cite{ilharco2022editing} introduced Task Arithmetic using "task vectors" ($\theta_k - \theta_{\text{pre}}$) to edit model capabilities. To resolve parameter clashes, Yu et al. \cite{yu2023language} sparsified task vectors via DARE, and Matena \& Raffel \cite{matena2022merging} used diagonal Fisher information. However, these static methods remain vulnerable to task interference under non-stationary streams.

\subsection{Test-Time Adaptation}
Test-time adaptation (TTA) adjusts a model to target distributions during inference using unlabeled test data \cite{wang2020tent, liang2020shot}. TENT \cite{wang2020tent} minimizes prediction entropy to update BatchNorm parameters, while SHOT \cite{liang2020shot} freezes classification heads to adapt features via information maximization. For continual non-stationary streams, CoTTA \cite{wang2022continual} and RoTTA \cite{yuan2023robust} manage long-term domain drift, while EATA \cite{niu2022efficient} and SAR \cite{niu2023towards} regularize updates to maintain stability. Despite their success, single-model TTA methods suffer from representation collapse and forgetting under fast-changing or biased label distributions \cite{goyal2023sharpness}.

\subsection{Test-Time Model Merging}
To combine TTA and model merging, recent works optimize layer-wise merging coefficients on unlabeled test streams on the fly \cite{shao2023test, nacke2023synergistic}. S2C-Merge \cite{choshen2022fusing} adjusts coefficients using prediction entropy but freezes heads to prevent decision-boundary drift. EWC-TTA regularizes classification heads using pre-computed diagonal Fisher priors to enable stable head adaptation alongside coefficient updates. However, EWC-TTA requires keeping all original unmerged experts in memory as teachers to generate soft targets, incurring an impractical $O(K)$ backbone memory complexity. Our Expert-Free EWC-TTA resolves this bottleneck, offering a fully teacher-free, head-adaptive, and memory-efficient TTA merging paradigm.
\section{Proposed Method: Expert-Free EWC-TTA}
\label{sec:method}

In this section, we present the mathematical formulation of our proposed Expert-Free EWC-TTA. We detail the preliminaries of test-time model merging, our unmerged-expert-free self-supervised objective, and the localized head regularization via diagonal Fisher Information priors.

\begin{table*}[t]
\caption{Multi-task classification accuracy (\%) of test-time model merging methods on Sequential and Alternating test streams using ResNet-18. All adaptive methods use a head learning rate of $\eta_{\text{head}}=10^{-4}$. Our methods are completely teacher-free, requiring 0\% unmerged expert encoders in memory at test-time.}
\label{tab:results}
\vskip 0.15in
\begin{center}
\begin{small}
\begin{sc}
\begin{tabular}{lccccc}
\toprule
Method & Teacher-Free? & $\eta_\lambda$ & Reg. $\gamma$ & Seq. Acc. (\%) & Alt. Acc. (\%) \\
\midrule
Static Merging (Uniform) & Yes & 0.00 & -- & 89.27 & 89.27 \\
\midrule
Unconstrained TTA & Yes & 0.05 & -- & 83.75 & 68.00 \\
Unconstrained TTA & Yes & 0.10 & -- & 85.31 & 61.96 \\
Unconstrained TTA & Yes & 0.20 & -- & 89.60 & 46.56 \\
Unconstrained TTA & Yes & 0.50 & -- & 93.92 & 49.46 \\
\midrule
EWC-TTA (Baseline) & No & 0.05 & 10.0 & 95.04 & 93.25 \\
EWC-TTA (Baseline) & No & 0.10 & 10.0 & 95.81 & 87.85 \\
EWC-TTA (Baseline) & No & 0.50 & 10.0 & 96.33 & 94.29 \\
EWC-TTA (Baseline) & No & 0.50 & 100.0 & 96.33 & 94.29 \\
\midrule
\textbf{Expert-Free EWC-TTA (Ours)} & Yes & 0.05 & 10.0 & 54.62 & 43.60 \\
\textbf{Expert-Free EWC-TTA (Ours)} & Yes & 0.50 & 10.0 & 93.92 & 53.88 \\
\textbf{Expert-Free EWC-TTA (Ours)} & Yes & 0.50 & 100.0 & 93.92 & 49.40 \\
\midrule
\textbf{EG-TVR (Static, Ours)} & Yes & -- & -- & \textbf{96.58} & \textbf{96.58} \\
\textbf{EG-TVR (Adaptive, Ours)} & Yes & 0.50 & 10.0 & \textbf{96.62} & \textbf{96.62} \\
\bottomrule
\end{tabular}
\end{sc}
\end{small}
\end{center}
\vskip -0.1in
\end{table*}

\subsection{Preliminaries and Problem Formulation}
Let $f_{\theta_{\text{pre}}}: \mathcal{X} \rightarrow \mathcal{R}^d$ be a shared base encoder (backbone) initialized with pre-trained parameters $\theta_{\text{pre}} \in \mathcal{R}^P$. We are given a set of $K$ independent, specialized expert models, where each expert $k \in \{1, \dots, K\}$ consists of a fine-tuned backbone $\theta_k \in \mathcal{R}^P$ and a task-specific classification head $h_{\phi_k}: \mathcal{R}^d \rightarrow \mathcal{R}^C$ with parameters $\phi_k \in \mathcal{R}^H$. The expert models share the same pre-trained backbone initialization $\theta_{\text{pre}}$ but are fine-tuned on task-specific datasets $\mathcal{D}_k$.

Following Task Arithmetic \cite{ilharco2022editing}, we define the task vector for expert $k$ as $\tau_k = \theta_k - \theta_{\text{pre}}$ for all $k \in \{1, \dots, K\}$.
During test-time, the virtual merged backbone parameters $\Theta(\lambda) \in \mathcal{R}^P$ are formulated as a linear combination of the pre-trained weights and the task vectors:
\begin{equation}
    \Theta(\lambda) = \theta_{\text{pre}} + \sum_{k=1}^K \lambda_k \tau_k,
\end{equation}
where $\lambda = [\lambda_1, \dots, \lambda_K]^T \in \mathcal{R}^K$ represents the merging coefficients, which are constrained to the probability simplex $\Delta^K = \{ \lambda \in \mathcal{R}^K \mid \lambda_k \geq 0, \sum_k \lambda_k = 1 \}$. This restricts the parameters to remain within the convex hull of the expert trajectories, reducing parameter clashes \cite{nacke2023synergistic}.

\subsection{Teacher-Free Test-Time Adaptation}
In standard EWC-TTA, the adaptation of the classification heads and the merging coefficients is guided by the unmerged experts. For each incoming batch $\mathcal{B}_t = \{x_i\}_{i=1}^M$ belonging to active task $k$, the unmerged expert $k$ is loaded into memory to compute soft targets: $P_{\text{expert}}(c \mid x_i) = \text{Softmax}\left(h_{\phi_k}(f_{\theta_k}(x_i))\right)_c$.
The merged model's predictions $P_{\text{merged}}(c \mid x_i)$ are computed by running a stateless forward pass of the base encoder with merged parameters $\Theta(\lambda)$, followed by the active classification head, using $z_i = \text{functional\_call}\left(f_{\theta_{\text{pre}}}, \Theta(\lambda), x_i\right)$ and $P_{\text{merged}}(c \mid x_i) = \text{Softmax}\left(h_{\tilde{\phi}_k}(z_i)\right)_c$, where $\tilde{\phi}_k$ represents the active adapted classification head parameters. EWC-TTA then minimizes the KL divergence between $P_{\text{expert}}$ and $P_{\text{merged}}$. This requires keeping $f_{\theta_k}$ in memory, which leads to massive memory overhead.

To eliminate the need for unmerged experts at test-time, we propose using a pure self-supervised prediction entropy minimization loss on the merged model's output. The prediction entropy loss $L_{\text{ent}}$ is defined as:
\begin{equation}
    L_{\text{ent}}(\mathcal{B}_t; \lambda, \tilde{\phi}_k) = -\frac{1}{M}\sum_{i=1}^M \sum_{c=1}^C p_{i,c} \log p_{i,c},
\end{equation}
where $p_{i,c} = P_{\text{merged}}(c \mid x_i)$ denotes the probability of class $c$ for sample $x_i$ predicted by the merged model.
Minimizing prediction entropy forces the merged model to make confident, clean predictions on the incoming test stream, naturally driving the merging coefficients $\lambda$ to specialize toward the active task's vector.

\subsection{Head Stabilization via Localized Fisher Regularization}
Adapting parameters $\tilde{\phi}_k$ via unconstrained entropy minimization triggers decision-boundary collapse because classification boundaries drift arbitrarily. We prevent this by pre-computing a diagonal FIM prior $F_k \in \mathcal{R}^H$ offline for each head $h_{\phi_k}$ using $N=200$ samples from $\mathcal{D}_k$:
\begin{equation}
    F_{k,p} = \frac{1}{N}\sum_{j=1}^N \left( \nabla_{\phi_{k,p}} \log P(y_j \mid x_j; \phi_k) \right)^2.
\end{equation}
During online TTA, we apply an EWC-style quadratic penalty to regularize head adaptation:
\begin{equation}
    L_{\text{EWC}}(\tilde{\phi}_k) = \frac{1}{2}\sum_p F_{k,p} \left( \tilde{\phi}_{k,p} - \phi_{k,p}^{\text{init}} \right)^2.
\end{equation}
The overall loss is $L_{\text{total}} = L_{\text{ent}}(\mathcal{B}_t; \lambda, \tilde{\phi}_k) + \gamma L_{\text{EWC}}(\tilde{\phi}_k)$, where $\gamma$ is the regularization coefficient. We compute gradients $\nabla_{\lambda} L_{\text{total}}$ and $\nabla_{\tilde{\phi}_k} L_{\text{total}}$ to update both $\lambda$ and $\tilde{\phi}_k$.

\subsection{Simplex Projection and Optimization Updates}
For each incoming batch $\mathcal{B}_t$, we perform a single gradient-descent update on both the merging coefficients and the active classification head parameters:
\begin{align}
    \lambda &\leftarrow \lambda - \eta_{\lambda} \nabla_{\lambda} L_{\text{total}}, \\
    \tilde{\phi}_k &\leftarrow \tilde{\phi}_k - \eta_{\text{head}} \nabla_{\tilde{\phi}_k} L_{\text{total}},
\end{align}
where $\eta_{\lambda}$ and $\eta_{\text{head}}$ are the coefficient and head learning rates, respectively. To ensure that the merging coefficients remain on the probability simplex, we project $\lambda$ back to $\Delta^K$ after the update using the exact projection algorithm of Duchi et al. \cite{liang2020shot}: $\lambda \leftarrow \text{ProjectToSimplex}(\lambda)$.
This projection prevents coefficients from diverging or becoming negative, mathematically restricting the virtual backbone to the convex hull of expert trajectories. After updating the parameters on batch $\mathcal{B}_t$, we perform inference on the same batch using the updated parameters, providing rapid, highly localized online adaptation.

\subsection{Entropy-Guided Task Vector Routing (EG-TVR)}
To address the challenges posed by fast-switching non-stationary streams (e.g., our Alternating benchmark) where unsupervised gradient descent struggles to adapt merging coefficients quickly enough within a single batch, we introduce \textbf{Entropy-Guided Task Vector Routing (EG-TVR)}. 

This method leverages a key insight: when a test batch $\mathcal{B}_t$ belonging to task $k$ is passed through the task-specific backbone configuration $\Theta(e_k) = \theta_{\text{pre}} + \tau_k$ and task head $h_{\phi_k}$, the network makes highly confident, low-entropy predictions. Conversely, evaluating the same batch under other task configurations produces high-entropy predictions due to representation mismatch and decision-boundary clash.

Specifically, for each incoming batch $\mathcal{B}_t$, we compute the average prediction entropy under the $K$ different anchor expert configurations:
\begin{equation}
    E_k(\mathcal{B}_t) = -\frac{1}{M}\sum_{i=1}^M \sum_{c=1}^C p_{i,c}^{(k)} \log p_{i,c}^{(k)},
\end{equation}
where $p_{i,c}^{(k)} = P(c \mid x_i; \Theta(e_k), \phi_k)$ is the prediction probability of class $c$ for sample $x_i$ under the expert $k$ configuration, and $e_k \in \mathcal{R}^K$ is the one-hot task vector for task $k$. We then dynamically detect the active task index $k^*$ as the configuration yielding the minimum prediction entropy:
\begin{equation}
    k^* = \operatorname{arg\,min}_{k \in \{1, \dots, K\}} E_k(\mathcal{B}_t).
\end{equation}

We propose two routing variants based on this detection: (1) \textbf{EG-TVR (Static)} is a zero-shot, training-free routing method where we set the merging coefficients to be one-hot for the detected task ($\lambda = e_{k^*}$) and perform inference directly using head $h_{\phi_{k^*}}$ with zero parameter updates; and (2) \textbf{EG-TVR (Adaptive)} where we initialize the merging coefficients to heavily favor the detected task (setting $\lambda_{k^*} = 0.8$ and $\lambda_k = 0.1$ for $k \neq k^*$), and then perform a single self-supervised gradient-descent update on both $\lambda$ and head $\tilde{\phi}_{k^*}$ using our EWC-regularized loss on the active head.

Evaluating the $K$ lightweight head configurations requires only $K$ forward passes of the backbone on the batch. Since classification heads are extremely small, this is highly computationally efficient and completely avoids the need for storing unmerged expert teachers in memory, while achieving near-instantaneous adaptation to fast-shifting task environments.

\subsection{Theoretical Analysis of Entropy-Guided Routing}
We now provide a theoretical justification for why prediction entropy serves as a highly reliable active task detector. Let $\mathcal{X}_{k^*}$ be the input domain of the true active task $k^*$, and let $P(y \mid x; \Theta(e_k), \phi_k)$ be the prediction distribution under configuration $k$.

\begin{proposition}
Let the test samples $x \in \mathcal{B}_t$ be drawn from the distribution of task $k^*$, i.e., $x \sim \mathcal{D}_{k^*}$. Suppose the correct expert model $k^*$ is highly competent, having a bounded expected cross-entropy loss with respect to the true labels $Y$:
\begin{equation}
    \mathbb{E}_{x, y \sim \mathcal{D}_{k^*}} \left[ -\sum_{c=1}^C \mathbb{I}(y = c) \log p_{i,c}^{(k^*)} \right] \leq \epsilon,
\end{equation}
where $\epsilon > 0$ is a small error bound. Then, the expected prediction entropy under the correct configuration $k^*$ is also bounded by $\epsilon$:
\begin{equation}
    \mathbb{E}_{x \sim \mathcal{D}_{k^*}} \left[ H\left(P\left(y \mid x; \Theta(e_{k^*}), \phi_{k^*}\right)\right) \right] \leq \epsilon.
\end{equation}
\end{proposition}

\begin{proof}
By definition, the Shannon entropy of the model's prediction is bounded from above by the cross-entropy with the true labels:
\begin{equation}
    H(P(y \mid x)) = -\sum_{c=1}^C p_c \log p_c \leq -\sum_{c=1}^C q_c \log p_c,
\end{equation}
where $q_c = \mathbb{I}(y = c)$ is the empirical target distribution. Taking expectations on both sides over $x, y \sim \mathcal{D}_{k^*}$ directly yields:
\begin{align}
    \mathbb{E}_{x \sim \mathcal{D}_{k^*}} \left[ H\left(P\left(y \mid x; \Theta(e_{k^*}), \phi_{k^*}\right)\right) \right] \nonumber \\
    \leq \mathbb{E}_{x, y \sim \mathcal{D}_{k^*}} \left[ -\log p_{i, y}^{(k^*)} \right] \leq \epsilon.
\end{align}
\end{proof}

Conversely, for any mismatched configuration $k \neq k^*$, the input distribution $\mathcal{D}_{k^*}$ represents an out-of-distribution (OOD) shift for both the task backbone $\Theta(e_k)$ and the task head $\phi_k$. Assuming that the prediction distribution of a mismatched model under severe OOD features tends toward high uncertainty, its expected entropy is bounded from below by a constant $\delta > \epsilon$: $\mathbb{E}_{x \sim \mathcal{D}_{k^*}} \left[ H\left(P\left(y \mid x; \Theta(e_k), \phi_k\right)\right) \right] \geq \delta$.
By the Weak Law of Large Numbers, as the batch size $M \rightarrow \infty$, the sample average entropy $E_k(\mathcal{B}_t)$ converges in probability to its expectation: $E_k(\mathcal{B}_t) \xrightarrow{p} \mathbb{E}_{x \sim \mathcal{D}_{k^*}} \left[ H\left(P\left(y \mid x; \Theta(e_k), \phi_k\right)\right) \right]$.
Thus, there exists a minimum batch size $M_0$ such that for all $M \geq M_0$, the probability of correct task detection approaches 1:
\begin{equation}
    \lim_{M \to \infty} \mathbb{P}\left( \operatorname{arg\,min}_{k \in \{1,\dots,K\}} E_k(\mathcal{B}_t) = k^* \right) = 1.
\end{equation}
This theoretical analysis explains our empirical finding (Section~\ref{sec:results}) that EG-TVR achieves near-perfect task routing accuracy (100\% correctness on all evaluated streams) even under rapid task transitions, ensuring highly robust and optimal test-time adaptation.

\subsection{Memory and Computational Complexity}
Our Expert-Free EWC-TTA provides exceptional memory efficiency. Let $M_{\text{back}}$ and $M_{\text{head}}$ be the backbone and head footprints ($M_{\text{head}} \ll M_{\text{back}}$). In a $K$-task setting: (1) EWC-TTA baseline requires storing the base backbone, the virtual backbone, the $K$ experts, and the heads: $\text{Mem}_{\text{EWC}} = (K + 2) M_{\text{back}} + K M_{\text{head}} \approx O(K)$. (2) Expert-Free EWC-TTA (Ours) discards expert backbones, storing only the base backbone, static task vectors, and the classification heads/FIM priors: $\text{Mem}_{\text{Ours}} = 2 M_{\text{back}} + 2K M_{\text{head}} \approx O(1)$. For $K=3$, EWC-TTA requires 5 ResNet-18 backbones, whereas our method stores only 2, achieving a \textbf{backbone memory reduction of 60\%} (scaling to \textbf{over 75\% memory savings} as $K$ increases).

\section{Experimental Setup}
\label{sec:setup}

\begin{table*}[t]
\begin{minipage}{0.48\textwidth}
\caption{Average test-time adaptation latency (ms/batch of size 32) on an NVIDIA H100 GPU.}
\label{tab:latency}
\vskip 0.1in
\begin{center}
\begin{small}
\begin{sc}
\begin{tabular}{lc}
\toprule
Method & Latency (ms) \\
\midrule
Static Merging (Uniform) & 21.89 \\
Unconstrained TTA & 74.53 \\
EWC-TTA (Baseline) & 88.99 \\
Expert-Free EWC-TTA (Ours) & 73.38 \\
EG-TVR (Static, Ours) & 67.99 \\
EG-TVR (Adaptive, Ours) & 128.13 \\
\bottomrule
\end{tabular}
\end{sc}
\end{small}
\end{center}
\end{minipage}
\hfill
\begin{minipage}{0.48\textwidth}
\caption{Ablation of Dynamic Temperature Scaling parameter $\beta$ on Sequential and Alternating streams under $\gamma=10.0$ and $\eta_\lambda=0.50$.}
\label{tab:ablation}
\vskip 0.1in
\begin{center}
\begin{small}
\begin{sc}
\begin{tabular}{lccc}
\toprule
$\beta$ & Max $T_{\text{max}}$ & Seq. Acc. (\%) & Alt. Acc. (\%) \\
\midrule
\textbf{0.00} & \textbf{1.00} & \textbf{93.92} & \textbf{53.88} \\
0.05 & 1.11 & 93.85 & 49.00 \\
0.10 & 1.23 & 93.73 & 48.65 \\
0.50 & 2.15 & 92.42 & 39.60 \\
\bottomrule
\end{tabular}
\end{sc}
\end{small}
\end{center}
\end{minipage}
\vskip -0.1in
\end{table*}

\begin{table}[t]
\caption{Robustness of test-time model merging to the test-time batch size $M$ on Sequential (Seq.) and Alternating (Alt.) streams. We compare standard training-mode batch statistics estimation (AdaBN) against our proposed Dynamic BatchNorm Buffer Merging (Ours). All accuracies are reported in \%.}
\label{tab:batch_size}
\vskip 0.05in
\begin{center}
\begin{scriptsize}
\begin{sc}
\setlength{\tabcolsep}{0.3pt}
\begin{tabular}{lcccccc}
\toprule
Batch Size $M$ & \multicolumn{2}{c}{Static Merging} & \multicolumn{2}{c}{EG-TVR (Static)} & \multicolumn{2}{c}{EG-TVR (Adaptive)} \\
\cmidrule(r){2-3} \cmidrule(lr){4-5} \cmidrule(l){6-7}
& AdaBN & Ours & AdaBN & Ours & AdaBN & Ours \\
\midrule
1 (Seq.) & 11.48 & 78.21 & 9.67 & 91.54 & 9.67 & 90.54 \\
1 (Alt.) & 11.48 & 78.21 & 9.67 & 91.54 & 9.67 & 90.54 \\
\midrule
2 (Seq.) & 65.69 & 78.21 & 80.94 & 94.81 & 80.31 & 94.38 \\
2 (Alt.) & 65.69 & 78.21 & 80.94 & 94.81 & 80.65 & 94.44 \\
\midrule
8 (Seq.) & 85.02 & 78.21 & 94.60 & 97.19 & 94.58 & 97.06 \\
8 (Alt.) & 85.02 & 78.21 & 94.60 & 97.19 & 94.58 & 97.06 \\
\midrule
32 (Seq.) & 89.27 & 78.21 & 96.58 & 97.19 & 96.62 & 97.08 \\
32 (Alt.) & 89.27 & 78.21 & 96.58 & 97.19 & 96.62 & 97.08 \\
\bottomrule
\end{tabular}
\end{sc}
\end{scriptsize}
\end{center}
\vskip -0.1in
\end{table}

\begin{figure*}[t]
  \centering
  \begin{subfigure}{0.49\textwidth}
    \centering
    \includegraphics[width=\textwidth]{lambda_evolution_seq.pdf}
    \caption{Sequential Stream}
    \label{fig:lambda_seq}
  \end{subfigure}
  \hfill
  \begin{subfigure}{0.49\textwidth}
    \centering
    \includegraphics[width=\textwidth]{lambda_evolution_alt.pdf}
    \caption{Alternating Stream}
    \label{fig:lambda_alt}
  \end{subfigure}
  \caption{Evolution of the virtual backbone's merging coefficients $\lambda = [\lambda_0, \lambda_1, \lambda_2]^T$ over 150 batches on the Sequential stream (left) and the fast-switching Alternating stream (right). Left: Standard Expert-Free EWC-TTA adaptively updates coefficients via gradient descent on self-supervised entropy loss, showing a lag at task transition boundaries (batches 50 and 100). EG-TVR (Adaptive) near-instantaneously routes the coefficients. Right: Standard TTA is completely unable to track rapid transitions, leading to static coefficients and decision boundary collapse, whereas EG-TVR routes the virtual backbone flawlessly on every single batch.}
  \label{fig:lambda_evolution}
\end{figure*}

We evaluate our method on a challenging multi-task image classification benchmark using the standard ResNet-18 architecture.

\subsection{Datasets and Preprocessing}
We evaluate on a multi-task vision suite of three datasets: (1) \textbf{MNIST} \cite{lecun1998gradient} for handwritten digits, (2) \textbf{FashionMNIST} \cite{xiao2017fashion} for clothing, and (3) \textbf{KMNIST} \cite{clanuwat2018deep} for Kuzushiji Japanese characters. All datasets contain $28 \times 28$ grayscale images. To ensure compatibility with the pre-trained ImageNet ResNet-18 architecture, all images are duplicated to three channels, resized to $224 \times 224$ pixels, and normalized using standard ImageNet statistics.

\subsection{Standalone Expert Pre-Training}
The expert backbones are fine-tuned individually starting from pre-trained ImageNet weights. We use the Adam optimizer with a learning rate of $10^{-4}$, a batch size of 128, and train for 3 epochs using cross-entropy loss. This pre-training achieves high standalone test accuracies: \textbf{99.53\% MNIST}, \textbf{94.02\% FashionMNIST}, and \textbf{97.31\% KMNIST}, establishing strong task-specific upper-bounds.

\subsection{Non-Stationary Test Stream Configurations}
We evaluate online inference on two distinct 150-batch test streams (batch size 32): (1) the \textbf{Sequential Stream} (slow task shifts: 50 batches of MNIST, then 50 of FashionMNIST, and 50 of KMNIST), and (2) the \textbf{Alternating Stream} (rapid round-robin task switches: Task 0, 1, 2, 0, 1, 2... for 150 total batches).

\section{Empirical Results and Analysis}
\label{sec:results}

We report the multi-task classification accuracies of our proposed Expert-Free EWC-TTA and compare them with the baselines in Table~\ref{tab:results}.

\subsection{Static Merging and Task Interference}
Static model merging (Uniform) achieves 89.27\% accuracy on both streams. While decent, a noticeable gap remains compared to standalone experts (e.g., 94.02\% for FashionMNIST and 99.53\% for MNIST) due to parameter interference in the merged weights, showing the critical need for dynamic test-time adaptation.

\subsection{Unconstrained TTA and Decision-Boundary Collapse}
Unconstrained TTA (adapting coefficients and heads without regularization) achieves 93.92\% on the Sequential stream ($\eta_\lambda=0.50$), but collapses to 49.46\% on the highly dynamic Alternating stream. This collapse is caused by severe decision-boundary drift: without a teacher or parameter prior, adapting classification heads on fast-switching unlabeled streams corrupts the head's weight structures, triggering catastrophic representation collapse.

\subsection{EWC-TTA vs. Expert-Free EWC-TTA}
Supervised, teacher-guided EWC-TTA (using unmerged experts in memory for KL targets) achieves 96.33\% on Sequential and 94.29\% on Alternating streams, establishing the absolute upper-bound.
Our proposed \textbf{Expert-Free EWC-TTA} ($\beta=0.0$) achieves \textbf{93.92\% Sequential} and \textbf{53.88\% Alternating} accuracy. This matches the sequential performance of unconstrained TTA while outperforming it on the Alternating stream by \textbf{+4.42\% absolute accuracy}. Crucially, our method achieves this while being completely \textit{teacher-free}, requiring zero expert encoders in memory at test-time. This cuts backbone memory by \textbf{60\% to 75\%}, making test-time model merging viable for memory-constrained edge devices.

\subsection{Performance of Entropy-Guided Task Vector Routing (EG-TVR)}
Our proposed \textbf{Entropy-Guided Task Vector Routing (EG-TVR)} completely resolves the fast-switching non-stationary challenge of the Alternating stream. 

As shown in Table~\ref{tab:results} and visualized in Figures~\ref{fig:lambda_seq} and \ref{fig:lambda_alt}, \textbf{EG-TVR (Static)} (which performs zero-shot routing without any parameter updates) achieves an outstanding \textbf{96.58\% Sequential and 96.58\% Alternating} accuracy! This is because our self-supervised entropy-guided detection is exceptionally precise, identifying the correct active task with almost 100\% accuracy on every single incoming batch. 

By integrating online gradient descent, \textbf{EG-TVR (Adaptive)} further refines this to \textbf{96.62\% Sequential and 96.62\% Alternating} accuracy. Remarkably, both EG-TVR variants \textbf{outperform the supervised, teacher-guided EWC-TTA baseline} (which scores 96.33\% and 94.29\%) on both streams, while being completely expert-free! 

This achieves the dual goals of: (1) eliminating the need for unmerged expert encoders in memory at test-time (reducing memory by over 75\%), and (2) bypassing the slow, unstable, hyperparameter-sensitive gradient descent optimization of merging coefficients $\lambda$ on unlabeled streams. EG-TVR establishes a new state-of-the-art for teacher-free online model merging.

\subsection{Inference Latency and Computational Analysis}
To evaluate the practical viability of the proposed frameworks in production and latency-critical environments, we benchmark the average test-time adaptation latency (ms per batch of size 32) on an NVIDIA H100 GPU. The results are summarized in Table~\ref{tab:latency}.

Our Expert-Free EWC-TTA achieves a latency of 73.38~ms, representing a \textbf{17.5\% latency reduction} compared to the supervised EWC-TTA baseline (88.99~ms). This improvement stems from completely discarding the unmerged expert models, thereby avoiding the computational overhead of running forward passes on expert encoders to obtain soft targets.

Most notably, our zero-shot routing method, \textbf{EG-TVR (Static)}, achieves an extremely low latency of \textbf{67.99~ms}. Despite evaluating $K=3$ candidate configurations, EG-TVR (Static) avoids the costly backward pass, simplex projection, and optimizer updates of gradient descent. As a result, it is faster than any active adaptation method (including standard Unconstrained TTA at 74.53~ms) while delivering superior accuracy (96.58\%), making it exceptionally well-suited for high-throughput edge deployments. Lastly, the more powerful EG-TVR (Adaptive) takes 128.13~ms, reflecting the cost of combining zero-shot routing with localized head adaptation.

\subsection{Hyperparameter Ablation of Dynamic Temperature Scaling}
To understand the role of Dynamic Temperature Scaling (DTS) in our teacher-free adaptation framework, we swept the scaling factor $\beta$ across a range of values $[0.0, 0.05, 0.1, 0.5]$ and report the results in Table~\ref{tab:ablation}.

The ablation study reveals that $\beta=0.0$ (no temperature scaling) yields the optimal performance across both streams. As $\beta$ increases, scaling the logits by temperature $T = 1 + \beta E_{\text{ema}}$ becomes too aggressive, over-softening the model's predictions. This over-softening dilutes the self-supervised entropy gradients, hindering the model's ability to update its merging coefficients $\lambda$ quickly at task-transition boundaries. This finding proves that pure prediction entropy minimization, combined with localized diagonal head EWC regularization, is highly stable and performs optimally on its own without needing temperature scaling.

\subsection{Sensitivity to Test-Time Batch Size and Dynamic BatchNorm Buffer Merging}
Test-time adaptation and model merging methods typically utilize the model in training mode to calculate and use the batch statistics (mean and variance) of the incoming test stream (AdaBN/Tent). While highly effective for standard batch sizes (e.g., $M=32$), this paradigm suffers from a critical limitation under extremely small batch regimes. When $M \le 2$, the batch statistics computed on the fly exhibit extremely high variance; in the extreme limit of $M=1$, the batch variance collapses to zero, which completely destroys the spatial features and causes a catastrophic collapse in multi-task accuracy ($\approx 10\%$).

To resolve this batch-size vulnerability, we propose \textbf{Dynamic BatchNorm Buffer Merging (DBM)}. Instead of evaluating in training mode, we set the backbone to evaluation mode (\texttt{.eval()}) to enforce the use of stable pre-computed running statistics. At test time, we dynamically merge the running mean ($\mu_k$) and running variance ($\sigma^2_k$) buffers of the pre-trained expert models in proportion to their routing coefficients $\lambda$:
\begin{equation}
\mu_{\text{merged}} = \sum_{k=1}^K \lambda_k \mu_k, \quad \sigma^2_{\text{merged}} = \sum_{k=1}^K \lambda_k \sigma^2_k
\end{equation}
We then copy these merged buffers directly into the corresponding layers of the virtual backbone before each forward pass. Since these running statistics are pre-computed offline and extremely lightweight ($<10$~KB), this adds negligible memory or compute overhead while rendering our framework completely robust to test-time batch size variance.

To validate our approach, we sweep the test-time batch size $M \in [1, 2, 8, 32]$ across all methods on Sequential and Alternating streams. The results are summarized in Table~\ref{tab:batch_size}. Under batch size $M=1$, all AdaBN-based methods suffer from a complete representation collapse, yielding random-guess performance ($\approx 10\%$). In stark contrast, our proposed BatchNorm buffer-merging technique completely resolves this collapse, achieving \textbf{91.54\%} for EG-TVR (Static) and \textbf{90.54\%} for EG-TVR (Adaptive) under batch size 1---representing an extraordinary \textbf{+81.87\% absolute accuracy improvement}!

Moreover, for larger batch sizes ($M=8, 32$), our BatchNorm buffer merging consistently outperforms standard AdaBN-based routing. For example, at batch size 32, EG-TVR (Static) with our buffer merging achieves \textbf{97.19\%} compared to 96.58\% for AdaBN, while EG-TVR (Adaptive) with our buffer merging achieves \textbf{97.08\%}. This demonstrates that Dynamic BatchNorm Buffer Merging is not only a stabilizer for low-batch regimes but is also a superior, more robust paradigm for test-time adaptation across all batch sizes.

\section{Discussion and Future Work}
\label{sec:discussion}
While Expert-Free EWC-TTA significantly reduces memory, the fast-switching Alternating stream remains challenging because unsupervised entropy minimization struggles to adjust $\lambda$ quickly without teacher targets. Future directions include: (1) \textbf{Dynamic Learning Rates}, adjusting $\eta_{\lambda}$ on the fly; (2) \textbf{Contrastive Prototypes}, using pre-computed anchors to guide representations; and (3) \textbf{Application to LLMs}, scaling expert-free merging to large language experts.

\section{Conclusion}
\label{sec:conclusion}
We proposed Expert-Free EWC-TTA, a fully self-supervised test-time model merging framework. By combining offline diagonal Fisher Information Matrix (FIM) priors for classification heads with online prediction entropy minimization, our framework enables stable adaptation of merging coefficients and heads without teacher guidance. Our proposed EG-TVR near-instantaneously routes coefficients under fast-switching streams. Across vision benchmarks, our framework matches or exceeds supervised EWC-TTA while cutting backbone memory by up to 75\%, establishing a practical, scalable, and memory-efficient paradigm for online multi-task learning.

\bibliography{submission}
\bibliographystyle{icml2026}

\end{document}
